\begin{table}[H]
\centering
\small
\begin{tabular}{llrrrr}
\toprule
Benchmark & $\lambda$ & $J^\lambda(\theta_\lambda^\star)$ & $|J^\lambda(\theta_\lambda^\star)-J(\theta^\star)|$ & $/\lambda$ & $/\lambda^2$ \\
\midrule
Linear--quadratic & $0.05$ & $7.499$ & $0.2754$ & $5.51$ & $110.17$ \\
 & $0.1$ & $8.325$ & $1.1015$ & $11.01$ & $110.15$ \\
 & $0.2$ & $11.628$ & $4.4047$ & $22.02$ & $110.12$ \\
 & $0.4$ & $24.836$ & $17.6126$ & $44.03$ & $110.08$ \\
 & $0.8$ & $77.654$ & $70.4299$ & $88.04$ & $110.05$ \\
\midrule
Portfolio & $0.025$ & $-13.164$ & $0.0308$ & $1.23$ & $49.34$ \\
 & $0.05$ & $-13.256$ & $0.1231$ & $2.46$ & $49.23$ \\
 & $0.1$ & $-13.622$ & $0.4887$ & $4.89$ & $48.87$ \\
 & $0.2$ & $-15.050$ & $1.9171$ & $9.59$ & $47.93$ \\
 & $0.4$ & $-20.479$ & $7.3465$ & $18.37$ & $45.92$ \\
\bottomrule
\end{tabular}
\caption{Perturbation bias of the optimal value on the two benchmarks where the perturbed objective and its optimizer are both available in closed form. Theorem~\ref{theorem:discrete-convergence-objective} bounds this quantity by $C_T\lambda$ in general, so the ratio in the fifth column is bounded and the sixth column would diverge if the rate were exactly linear. It does not: the sixth column is constant to within $0.1\%$ over a sixteen-fold range of $\lambda$ on the linear--quadratic benchmark, and constant to within $7\%$ on the portfolio benchmark. Both benchmarks are quadratic in the population argument, for which the linear term of the expansion vanishes and the bias is $O(\lambda^2)$; the measurement recovers that rate rather than the generic bound.}
\label{tab:objective-bias-scaling}
\end{table}
